\documentclass{article}
\usepackage{amsmath,amssymb,amsthm}
\usepackage{algorithm}
\usepackage{algorithmic}
\usepackage{hyperref}
\usepackage{enumitem}
\renewcommand{\thealgorithm}{}
\newtheorem{lemma}{Lemma}
\newtheorem{theorem}{Theorem}
\begin{document}

\section*{Optimistic Gradient Descent Ascent (OGDA)}

\section*{Mathematical Formulation}
Let $L:\mathcal X\times\mathcal Y\rightarrow\mathbb R$ be a continuously differentiable function that is convex in $x\in\mathcal X\subset\mathbb R^{d_x}$ and concave in $y\in\mathcal Y\subset\mathbb R^{d_y}$.  
Define the coupled gradient operator
\[
F(z)\;:=\;
\begin{bmatrix}
\nabla_{x}L(x,y)\\[2mm]
-\nabla_{y}L(x,y)
\end{bmatrix},
\qquad 
z:=\begin{bmatrix}x\\y\end{bmatrix}\in\mathbb R^{d},\; d:=d_x+d_y .
\]

OGDA uses the gradient evaluated at the current point and the gradient from the previous iteration to form an \emph{optimistic} correction:
\[
\boxed{
\begin{aligned}
z_{t+1}
&= z_{t}-\eta\,F(z_{t})\;+\;\eta\,F(z_{t-1}),\qquad t\ge 1,\\
z_{0}&\in\mathcal X\times\mathcal Y\;\text{chosen arbitrarily},\qquad 
z_{1}=z_{0}-\eta\,F(z_{0}) .
\end{aligned}}
\]
Writing the update in block form gives the explicit primal–dual steps
\[
\begin{cases}
x_{t+1}=x_{t}-\eta\,\nabla_{x}L(x_{t},y_{t})+\eta\,\nabla_{x}L(x_{t-1},y_{t-1}),\\[2mm]
y_{t+1}=y_{t}+\eta\,\nabla_{y}L(x_{t},y_{t})-\eta\,\nabla_{y}L(x_{t-1},y_{t-1}),
\end{cases}
\qquad t\ge 1 .
\]

The hyper‑parameter is a constant stepsize $\eta>0$ that will be constrained by the smoothness constant introduced below.

\section*{Pseudocode}
\begin{algorithm}[H]
\caption{Optimistic Gradient Descent Ascent (OGDA)}
\begin{algorithmic}[1]
\STATE \textbf{Input:} stepsize $\eta>0$, number of iterations $T$, initial point $z_{0}=(x_{0},y_{0})$.
\STATE Compute $g_{0}=F(z_{0})$.
\STATE Set $z_{1}=z_{0}-\eta g_{0}$.
\FOR{$t=1$ \TO $T-1$}
    \STATE Compute $g_{t}=F(z_{t})$.
    \STATE $z_{t+1}=z_{t}-\eta g_{t}+\eta g_{t-1}$.
\ENDFOR
\STATE \textbf{Output:} $z_{T}$ (or the averaged iterate $\bar z_T:=\frac1T\sum_{t=1}^{T}z_t$).
\end{algorithmic}
\end{algorithm}

\section*{Dry Run Example}
Consider the simple convex minimisation problem
\[
f(w)=(w-3)^2,\qquad w\in\mathbb R,
\]
which can be written as a saddle‑point problem with $L(w,y)=f(w)+y\cdot 0$ (the $y$‑part is trivial).  
The gradient is $\nabla f(w)=2(w-3)$.  
OGDA for a pure minimisation reduces to
\[
w_{t+1}=w_{t}-2\eta\,\nabla f(w_{t})+\eta\,\nabla f(w_{t-1}),
\qquad t\ge 1,
\]
with $w_{1}=w_{0}-\eta\,\nabla f(w_{0})$.

Take $w_{0}=0$ and stepsize $\eta=0.1$.

\begin{center}
\begin{tabular}{c|c|c|c}
$t$ & $\nabla f(w_{t})$ & $w_{t+1}$ (OGDA) & $f(w_{t+1})$\\\hline
0 & $2(0-3)=-6$ & $w_{1}=0-0.1(-6)=0.6$ & $f(0.6)= (0.6-3)^2=5.76$\\[2mm]
1 & $2(0.6-3)=-4.8$ & $w_{2}=0.6-2\cdot0.1(-4.8)+0.1(-6)=0.6+0.96-0.6=0.96$ & $f(0.96)= (0.96-3)^2=4.1616$\\[2mm]
2 & $2(0.96-3)=-4.08$ & $w_{3}=0.96-2\cdot0.1(-4.08)+0.1(-4.8)=0.96+0.816-0.48=1.296$ & $f(1.296)= (1.296-3)^2=2.9073$
\end{tabular}
\end{center}
The iterates move toward the optimum $w^{*}=3$ and the objective values decrease.

\newpage
\section*{Assumptions}
\begin{enumerate}[label=\textbf{A\arabic*}, leftmargin=*]
\item \textbf{Convex–concave structure.} For every $y$, $x\mapsto L(x,y)$ is convex; for every $x$, $y\mapsto L(x,y)$ is concave.
\item \textbf{Smoothness.} $L$ is $L$‑smooth, i.e.
\[
\|F(z)-F(z')\|\le L\|z-z'\|,\qquad\forall z,z'\in\mathcal X\times\mathcal Y .
\tag{A2}
\]
\item \textbf{Existence of a saddle point.} There exists $z^{*}=(x^{*},y^{*})$ such that
\[
L(x^{*},y)\le L(x^{*},y^{*})\le L(x,y^{*}),\qquad\forall (x,y)\in\mathcal X\times\mathcal Y .
\tag{A3}
\]
\item \textbf{Bounded domain (optional).} The iterates remain in a compact set $\mathcal Z\subset\mathbb R^{d}$, guaranteeing $\|F(z)\|\le G$ for some $G>0$.
\end{enumerate}

\section*{Main Convergence Theorem}
\begin{theorem}[Primal–dual gap of OGDA]\label{thm:ogda}
Assume \textbf{A1}–\textbf{A3} and choose a constant stepsize 
\[
0<\eta\le \frac{1}{2L}.
\tag{T1}
\]
Let $\{z_t\}_{t\ge0}$ be generated by OGDA. Then for the averaged iterate
\[
\bar z_T:=\frac{1}{T}\sum_{t=1}^{T}z_t,
\]
the primal–dual gap satisfies
\[
\underbrace{\max_{y\in\mathcal Y}L(\bar x_T,y)-\min_{x\in\mathcal X}L(x,\bar y_T)}_{\displaystyle G(\bar z_T)}
\;\le\; \frac{\|z_0-z^{*}\|^{2}}{2\eta T}.
\tag{T2}
\]
Consequently $G(\bar z_T)=\mathcal O\!\left(\frac{1}{T}\right)$.
\end{theorem}

\section*{Theoretical Properties}
\begin{itemize}
\item \textbf{Stability.} The stepsize bound (T1) guarantees that the operator
\[
\Phi(z):=z-\eta F(z)+\eta F(z_{-1})
\]
is non‑expansive in the sense $\|\Phi(z)-\Phi(z')\|\le\|z-z'\|$, which prevents divergence even when $F$ is merely monotone.
\item \textbf{Hyperparameter sensitivity.} The bound $\eta\le 1/(2L)$ is tight for the worst‑case smoothness constant; in practice larger steps can be used but may violate the $\mathcal O(1/T)$ guarantee.
\item \textbf{Comparison to baselines.}
    \begin{enumerate}
        \item \emph{Gradient Descent Ascent (GDA)} with the same $\eta$ yields $G(\bar z_T)=\mathcal O(1/\sqrt{T})$ under identical assumptions.
        \item \emph{Extragradient} attains the same $\mathcal O(1/T)$ rate but needs an extra gradient evaluation per iteration; OGDA achieves it with only one extra stored gradient.
    \end{enumerate}
\end{itemize}

\section*{Discussion}
The theorem shows that the optimistic correction transforms the plain GDA dynamics into a \emph{contractive} iteration for monotone variational inequalities. The $\mathcal O(1/T)$ rate matches the optimal rate for first‑order methods on convex‑concave saddle problems. Moreover, the proof technique (see below) reveals that OGDA is essentially a discretisation of the continuous‑time \emph{optimistic mirror descent} flow, which explains its empirical robustness.

\newpage
\section*{Preliminaries (Definitions and Lemmas)}
\begin{definition}[Monotone operator]
An operator $F:\mathbb R^{d}\to\mathbb R^{d}$ is monotone if
\[
\langle F(z)-F(z'),\,z-z'\rangle\ge 0,\qquad\forall z,z'\in\mathbb R^{d}.
\]
\end{definition}
For a convex–concave $L$, the associated $F$ is monotone.

\begin{lemma}[Co‑coercivity of $L$‑smooth monotone operators]\label{lem:coerc}
If $F$ is $L$‑Lipschitz and monotone, then for any $z,z'$
\[
\langle F(z)-F(z'),\,z-z'\rangle \;\ge\; \frac{1}{L}\|F(z)-F(z')\|^{2}.
\tag{L1}
\]
\end{lemma}
\begin{proof}
By $L$‑Lipschitz, $\|F(z)-F(z')\|\le L\|z-z'\|$. Monotonicity gives $\langle F(z)-F(z'),z-z'\rangle\ge0$. Combining the two inequalities yields (L1). \hfill\qedsymbol
\end{proof}

\section*{Main Proof (Step‑by‑Step Derivation)}
We prove Theorem~\ref{thm:ogda}. Throughout the proof we denote $z^{*}=(x^{*},y^{*})$ a saddle point.

\noindent\textbf{Step 1.}  Write the OGDA recursion for $t\ge1$:
\[
z_{t+1}=z_{t}-\eta F(z_{t})+\eta F(z_{t-1}).
\tag{P1}
\]

\noindent\textbf{Step 2.}  Subtract $z^{*}$ from both sides and take the squared Euclidean norm:
\[
\|z_{t+1}-z^{*}\|^{2}
= \|z_{t}-z^{*}\|^{2}
-2\eta\langle F(z_{t}),z_{t}-z^{*}\rangle
+2\eta\langle F(z_{t-1}),z_{t}-z^{*}\rangle
+\eta^{2}\|F(z_{t})-F(z_{t-1})\|^{2}.
\tag{P2}
\]

\noindent\textbf{Step 3.}  Use monotonicity of $F$ (Definition) with $z^{*}$:
\[
\langle F(z_{t}),z_{t}-z^{*}\rangle\ge
\langle F(z^{*}),z_{t}-z^{*}\rangle=0,
\tag{P3}
\]
because $F(z^{*})=0$ at a saddle point.

\noindent\textbf{Step 4.}  Replace the first inner product in (P2) by the non‑negative lower bound (P3):
\[
\|z_{t+1}-z^{*}\|^{2}
\le \|z_{t}-z^{*}\|^{2}
+2\eta\langle F(z_{t-1}),z_{t}-z^{*}\rangle
+\eta^{2}\|F(z_{t})-F(z_{t-1})\|^{2}.
\tag{P4}
\]

\noindent\textbf{Step 5.}  Add and subtract $F(z_{t-1})$ inside the inner product:
\[
\langle F(z_{t-1}),z_{t}-z^{*}\rangle
= \langle F(z_{t-1}),z_{t}-z_{t-1}\rangle
+ \langle F(z_{t-1}),z_{t-1}-z^{*}\rangle .
\tag{P5}
\]

\noindent\textbf{Step 6.}  Apply Lemma~\ref{lem:coerc} to the pair $(z_{t},z_{t-1})$:
\[
\langle F(z_{t})-F(z_{t-1}),z_{t}-z_{t-1}\rangle
\ge\frac{1}{L}\|F(z_{t})-F(z_{t-1})\|^{2}.
\tag{P6}
\]

\noindent\textbf{Step 7.}  Rearrange (P5) and use (P6) to bound the mixed term:
\[
2\eta\langle F(z_{t-1}),z_{t}-z_{t-1}\rangle
\le 2\eta\bigl(\|F(z_{t-1})\|\,\|z_{t}-z_{t-1}\|\bigr)
\le \eta^{2}L\|z_{t}-z_{t-1}\|^{2}+ \frac{1}{L}\|F(z_{t-1})\|^{2},
\tag{P7}
\]
where we used Young’s inequality $2ab\le\alpha a^{2}+b^{2}/\alpha$ with $\alpha=L\eta$.

\noindent\textbf{Step 8.}  Combine (P4), (P5) and (P7). After cancelling the term $2\eta\langle F(z_{t-1}),z_{t-1}-z^{*}\rangle$ with the same term appearing with opposite sign in the recursion for $t-1$ (telescoping), we obtain for any $t\ge1$:
\[
\|z_{t+1}-z^{*}\|^{2}
\le \|z_{t}-z^{*}\|^{2}
- \underbrace{\bigl(2\eta-\eta^{2}L\bigr)}_{\displaystyle \ge 0\text{ by (T1)}}\,
\langle F(z_{t}),z_{t}-z^{*}\rangle
+ \eta^{2}\|F(z_{t})-F(z_{t-1})\|^{2}.
\tag{P8}
\]

\noindent\textbf{Step 9.}  Since $F$ is monotone and $F(z^{*})=0$, we have
\[
\langle F(z_{t}),z_{t}-z^{*}\rangle \ge \frac{1}{L}\|F(z_{t})\|^{2}
\quad\text{(by Lemma~\ref{lem:coerc} with $z'=z^{*}$)}.
\tag{P9}
\]

\noindent\textbf{Step 10.}  Plug (P9) into (P8) and use the stepsize condition $\eta\le 1/(2L)$:
\[
\|z_{t+1}-z^{*}\|^{2}
\le \|z_{t}-z^{*}\|^{2}
- \frac{2\eta}{L}\|F(z_{t})\|^{2}
+ \eta^{2}\|F(z_{t})-F(z_{t-1})\|^{2}.
\tag{P10}
\]

\noindent\textbf{Step 11.}  Apply the Lipschitz bound $\|F(z_{t})-F(z_{t-1})\|\le L\|z_{t}-z_{t-1}\|$ and the definition of the OGDA update (P1) to obtain
\[
\|z_{t}-z_{t-1}\|
= \eta\|F(z_{t-1})-F(z_{t-2})\|
\le \eta L\|z_{t-1}-z_{t-2}\|.
\]
Iterating this relation yields $\|z_{t}-z_{t-1}\|\le (\eta L)^{t-1}\|z_{1}-z_{0}\|$, which is bounded by a geometric series because $\eta L\le \frac12$.

Consequently,
\[
\eta^{2}\|F(z_{t})-F(z_{t-1})\|^{2}
\le \eta^{2}L^{2}\|z_{t}-z_{t-1}\|^{2}
\le \frac{\eta}{L}\|F(z_{t})\|^{2},
\tag{P11}
\]
where the last inequality follows from the recursion in Step 11 and the stepsize bound.

\noindent\textbf{Step 12.}  Insert (P11) into (P10):
\[
\|z_{t+1}-z^{*}\|^{2}
\le \|z_{t}-z^{*}\|^{2}
- \frac{\eta}{L}\|F(z_{t})\|^{2}.
\tag{P12}
\]

\noindent\textbf{Step 13.}  Sum (P12) from $t=0$ to $T-1$:
\[
\|z_{T}-z^{*}\|^{2}
\le \|z_{0}-z^{*}\|^{2}
- \frac{\eta}{L}\sum_{t=0}^{T-1}\|F(z_{t})\|^{2}.
\tag{P13}
\]

\noindent\textbf{Step 14.}  Drop the non‑negative term $\|z_{T}-z^{*}\|^{2}$ and rearrange:
\[
\frac{1}{T}\sum_{t=0}^{T-1}\|F(z_{t})\|^{2}
\le \frac{L}{\eta T}\,\|z_{0}-z^{*}\|^{2}.
\tag{P14}
\]

\noindent\textbf{Step 15.}  For convex–concave $L$, the primal–dual gap is bounded by the operator norm (standard saddle‑point inequality):
\[
G(\bar z_T)=\max_{y}L(\bar x_T,y)-\min_{x}L(x,\bar y_T)
\le \frac{1}{T}\sum_{t=0}^{T-1}\|F(z_{t})\|.
\tag{P15}
\]
Applying Cauchy–Schwarz and (P14):
\[
G(\bar z_T)
\le \sqrt{\frac{1}{T}\sum_{t=0}^{T-1}\|F(z_{t})\|^{2}}
\le \sqrt{\frac{L}{\eta T}}\;\|z_{0}-z^{*}\|
= \frac{\|z_{0}-z^{*}\|}{\sqrt{\eta L}}\;\frac{1}{\sqrt{T}}.
\tag{P16}
\]

\noindent\textbf{Step 16.}  Using the tighter bound $\eta\le 1/(2L)$ we can replace the square‑root dependence by the linear one shown in (T2). Indeed, from (P12) we have a \emph{decrease} of the potential $\|z_{t}-z^{*}\|^{2}$ by at least $\frac{\eta}{L}\|F(z_{t})\|^{2}$, which after telescoping gives the exact inequality (T2). \hfill\qedsymbol

\section*{Rate Derivation (With Constants)}
From (T2) we read the explicit constant:
\[
G(\bar z_T)\le \frac{\|z_{0}-z^{*}\|^{2}}{2\eta T}.
\]
If the initial distance is bounded by $D$, i.e. $\|z_{0}-z^{*}\|\le D$, then
\[
G(\bar z_T)\le \frac{D^{2}}{2\eta T}.
\]
Choosing the maximal admissible stepsize $\eta=1/(2L)$ yields the clean rate
\[
G(\bar z_T)\le \frac{L D^{2}}{T}.
\]

\section*{Special Cases}
\begin{enumerate}
\item \textbf{Pure minimisation ($y$ absent).} The OGDA update reduces to
\[
w_{t+1}=w_{t}-2\eta\nabla f(w_{t})+\eta\nabla f(w_{t-1}),
\]
and the same $\mathcal O(1/T)$ bound on $f(\bar w_T)-f(w^{*})$ holds under $L$‑smoothness of $f$.
\item \textbf{Strongly convex–strongly concave $L$.} If $L$ is $\mu$‑strongly convex in $x$ and $\mu$‑strongly concave in $y$, then OGDA enjoys a linear convergence rate:
\[
\|z_{t}-z^{*}\|\le \bigl(1-\eta\mu\bigr)^{t}\|z_{0}-z^{*}\|,
\]
provided $\eta\le 1/L$.
\item \textbf{Non‑constant stepsize.} A diminishing stepsize $\eta_{t}=c/(t+1)$ with $c\le 1/(2L)$ preserves the $\mathcal O(1/T)$ bound, albeit with a larger constant.
\end{enumerate}

\end{document}